% C-49 -- Hunting Instinct naming/capitalization consistency (redundancy_closing_audit_consolidation_2026-09-02.md, P-20)
% Audit copy only -- live files are sections/methodology.tex, sections/results.tex.
%
% Confirmed by Samuel and Diego. $HI$ was introduced as "hunting drive" in methodology.tex:235 and
% "hunting-instinct" (adjective form) in methodology.tex:292 -- two names for the same parameter --
% while results.tex used "hunting instinct" (line 216, already lowercase) and "Hunting Instinct"
% (line 652, capitalized) -- not a literature term, an internal model name, no reason to capitalize.
% Unified on "hunting instinct" (noun) / "hunting-instinct" (adjective) everywhere, all lowercase.

% --- methodology.tex, ~line 235 ---
$HI$ denotes an intrinsic \revblue{hunting instinct} that provides a baseline forward/exploratory motivation, allowing the platform to keep moving and exploring even in the absence of external stimuli.

% --- results.tex, ~line 652 ---
\revblue{A simultaneous $X_{15}$/$X_{16}$ activation (Figure~\ref{fig:NEU_EST_MUL_real}d), from the constant \revblue{hunting-instinct} stimulus, keeps $X_{16}$ active and uninhibited by $X_{15}$ --- the no-motor-action state shown in Figure~\ref{fig:frames-c}.}
